% Source: https://artofproblemsolving.com/community/c3h1649579_problem_set_4
% Stable identity: aops-topic-1649579-q03
The number $m$ is the \emph{slope} of the line: the change in $y$ for each increase of $1$ in $x$, so between any two points of the line
\[
m=\dfrac{\text{rise}}{\text{run}}=\dfrac{y_2-y_1}{x_2-x_1}.
\]
(The number $b$ is the $y$-intercept: the line crosses the $y$-axis at $(0,b)$.)

For the points $(1,2)$ and $(3,8)$, $m=\dfrac{8-2}{3-1}=3$. Then $2=3\cdot1+b$ gives $b=-1$, so the line is $y=3x-1$:
\[
\boxed{m=3,\quad b=-1}.
\]

\begin{asy}
usepackage("amsmath");
size(4.5cm,7cm);
defaultpen(fontsize(10pt));
import graph;
draw((-1.5,0)--(3.8,0),gray,Arrow(TeXHead));draw((0,-2)--(0,9.5),gray,Arrow(TeXHead));
label("$x$",(3.8,0),E,gray);label("$y$",(0,9.5),N,gray);
draw((-0.33,-2)--(3.4,9.2),linewidth(0.9));
pair P=(1,2),Q=(3,8);
draw(P--(3,2),blue+dashed);draw((3,2)--Q,blue+dashed);
dot(P);dot(Q);dot((0,-1));
label("$(1,2)$",P,NW);label("$(3,8)$",Q,W);label("$(0,-1)$",(0,-1),W);
label("run $2$",(2,2),S,blue);label("rise $6$",(3,5),E,blue);
label("\mbox{}",(min(currentpicture,user=true)+max(currentpicture,user=true))/2,black); // keeps the text after the figure black when LaTeX typesets the labels inline
\end{asy}
